\magicSection{Managed Ring Buffers}{sec:ManagedRingBuffers}

In general, the mapping between device tasks (def~\ref{sec:gDeviceTask}) and CUDA clusters (def~\ref{sec:gCluster}) is unpredictable.
This poses difficulties for expressing ring buffer usage, and more generally, coordinating shared memory usage between subsequent device tasks assigned to the same CUDA cluster (persistent kernel).

Frequently, the last ring buffer slot used for a given device task may not be the last physical slot.
To minimize latency, the next device task will start off using the least-recently-used ring buffer slot, instead of physical slot 0.

For example, if a kernel uses 5 ring buffer slots, and the ring buffer has depth 3, then the first device task on a CUDA cluster will use ring buffer slots in sequence \texttt{[0, 1, 2, 0, 1]}.
The next device task on the cluster will use ring buffer slots in sequence \texttt{[2, 0, 1, 2, 0]}.

We cannot express this within the limitations of Exo's indexing restrictions, since the device task / cluster mapping relies on runtime values.
Exo-GPU handles this problem by introducing ``managed ring buffers'', where the modular arithmetic for the ring buffering is hidden in the implementation, rather than explicitly in the object language.
There is a hidden \lighttt{consumption} variable for each managed ring buffer that coordinates the handoff between subsequent users of the buffer.

An allocation is a managed ring buffer iff one dimension is annotated with \lighttt{ring\_buffer\_by}.
This is the \myKeyA{managed ring buffer dimension}.
Currently, these allocations must be CUDA SMEM or mbarrier allocations.

To make this safe, we ``guard'' each managed ring buffer using mbarriers (Section~\ref{sec:MbarrierUsage}).
Data allocations require an explicit \lighttt{.ring\_guarded\_by} annotation; barriers do not (they guard themselves).
The intent of this is to establish an implicit contract between the programmer and the underlying SMEM allocator.
The requirements of Section~\ref{sec:ManagedRingBufferSync} enforce that the program arrives on an mbarrier to relinquish usage of a certain ring buffer slot, then awaits on the same physical mbarrier to gain permission to use another ring buffer slot aliased to the one previously relinquished.

\mainKey{Scheduling Functions:}

\texttt{set\_dim\_size\_annotation(\codecomment{...}, dim\_idx, ring\_buffer\_by(\codecomment{...}))}

\texttt{set\_ring\_guarded\_by(\codecomment{...}, modify\_cursor, guarded\_by\_cursor)}

\texttt{insert\_barrier\_alloc(\codecomment{...}, CudaMbarrierPreArrive(\codecomment{...}), ring\_dim\_idx, ring\_depth)}

\magicSubsection{Managed Ring Buffer Indexing}{sec:ManagedRingBufferIndexing}

The \lighttt{consumption} value for a managed ring buffer is tracked per-CTA, and persists for the lifetime of the kernel.
It is initially 0, and not reset for each new device task.

Let \lighttt{ring\_logical\_extent @ ring\_buffer\_by(ring\_depth)} be the extent of the managed ring buffer dimension.
For data allocations, let \lighttt{pre\_arrive = 0}.
For mbarrier allocations, \lighttt{pre\_arrive} is the parameter to \lighttt{CudaMbarrierPreArrive}.
During codegen, we do the following rewrites:

\begin{itemize}
  \item Rewrite the physical extent of the managed ring buffer dimension to \lighttt{ring\_depth}.
  \item Rewrite values \lighttt{idx} that index the managed ring buffer dimension to \lighttt{offset\_idx \% ring\_depth},
    with \lighttt{offset\_idx = idx + consumption - pre\_arrive}.
  \item Insert \lighttt{consumption += ring\_logical\_extent - pre\_arrive} immediately after the free
  \item For mbarriers, the parity for the \lighttt{Await} is given by \lighttt{offset\_idx / ring\_depth \% 2}.
  \item Skip Arrive/Await for mbarriers where \lighttt{offset\_idx < 0}.
\end{itemize}

Extents and indices on other dimensions are not modified.


\magicSubsection{Managed Ring Buffer Sync}{sec:ManagedRingBufferSync}

Each element of a managed ring buffer has associated \myKeyA{guard barriers} and an optional \myKeyA{free-on-arrive} barrier.
These respectively describe the awaits and arrives expected.

\mainKey{SMEM Data Allocations:} Let $x[j^*]$ be an element of the managed ring buffer data allocation.
Let $z$ be the barrier variable named by the \lighttt{.ring\_guarded\_by($z$)} of the $x$ allocation.

With $Z$ being the dimensionality of $z$, let $j_z^*$ be the first $Z$ coordinates of $j^*$.
Let $\widehat{u}$ be a $Z$-dimensional unit vector with 1 on the managed ring buffer dimension.
This managed ring buffer dimension must agree between $x$ and $z$.
Then $x[j^*]$ has
\begin{itemize}
  \item guard barriers being the multicast group of $z[j_z^*]$ (def~\ref{sec:gMulticastGroup})
  \item free-on-arrive barrier $z[j_z^* + (\texttt{ring\_depth}) \widehat{u}]$.
\end{itemize}


For example, an allocation\\
\blacktt{x: dtype[A.ring\_buffer\_by(ring\_depth), B, C].ring\_guarded\_by(z)}\\
(with \blacktt{z: barrier[A + ring\_depth, B]})\\
will have for each \blacktt{x[a, b, c]}
\begin{itemize}
  \item guard barriers being the multicast group of \blacktt{z[a, b]} (if no multicasting, the group is \blacktt{z[a, b]} alone)
  \item free-on-arrive barrier \blacktt{z[a + ring\_depth, b]}.
\end{itemize}

There must be no out-of-range barrier array elements referenced in the above calculations.

\mainKey{mbarrier Allocations:} Let $z[j^*]$ be an element of a managed ring buffer barrier allocation.
With $Z$ being the dimensionality of $z$, let $\widehat{u}$ be a $Z$-dimensional unit vector with 1 on the managed ring buffer dimension.
Then $z[j^*]$ has
\begin{itemize}
  \item guard barriers being the multicast group of $z[j^* - (\texttt{ring\_depth}) \widehat{u}]$ (def~\ref{sec:gMulticastGroup})
  \item free-on-arrive barrier $z[j^* + (\texttt{ring\_depth}) \widehat{u}]$.
\end{itemize}

Out-of-range array elements are removed in the above calculations.
This means the last few barrier array elements may have no free-on-arrive barrier.
\redBox{Note:} this introduces a blind spot where mbarrier usage may be unsafe due to interactions across different device tasks, due to usage patterns that are not problematic when analyzing device tasks in isolation.

\mainKey{Await Requirement:} Must await for at least one of the guard barriers of $x[j^*]$ before any reads, mutates, or sync statement instances referencing $x[j^*]$, unless there are no guard barriers.
This is enforced in the abstract machine by adding a special \textsf{VisRecord} upon allocation, with an empty visibility set (Section~\ref{sec:SyncSemanticsManagedRingBuffer}).

\mainKey{Arrive Requirement:} If a free-on-arrive barrier exists for $x[j^*]$, the program must arrive on the free-on-arrive barrier after all reads, mutates, and sync statement instances referencing $x[j^*]$.
Executed statements must not reference $x[j^*]$ after the arrive.
This is enforced by the abstract machine semantics upon free (Section~\ref{sec:ChecksOnFree}).
The abstract machine checks for a mandatory pending await (def~\ref{sec:gPendingAwait}) referencing the free-on-arrive barrier, which will be missing for any reads or mutates that were not guaranteed to retire before the \lighttt{Arrive} in question.

